\documentclass{obelus-preprint}
\usepackage{verbatim}


\title{Testing LaTeX Markdown Compatibility}
\author{Heather Leffew, PhD\\Obelus Institute}
\date{May 2026}

\begin{document}
\maketitle

\begin{abstract}
This is a test abstract to verify that the \texttt{md\_to\_latex.py} script correctly parses the dual-header structure, successfully dropping the YAML frontmatter and capturing the formal title block.
\end{abstract}
\section{Introduction}
Here is the first section of the document. We are testing whether standard markdown features like \textbf{bolding}, \emph{italics}, and \texttt{inline code} work as expected.
\section{Methods}
We enforce a strict prohibition on raw HTML tags. Instead of an HTML table, we use a pure markdown table:
\begin{center}
\begin{tabular}{lll}
\toprule
\textbf{Metric} & \textbf{Value} & \textbf{Status} \\
\midrule
Test 1 & 95\% & Pass \\
Test 2 & 88\% & Pass \\
\bottomrule
\end{tabular}
\end{center}
There is also no raw HTML for visualizations. Any JS scripts are loaded via the frontmatter.
\section{Conclusion}
If the script parses this successfully, it means the new workflow guidelines are fully compatible with formal LaTeX preprint submissions.
\end{document}